\documentclass[11pt,a4paper]{article}
\usepackage[utf8]{inputenc}
\usepackage{amsmath,amssymb}
\usepackage{graphicx}
\usepackage{booktabs}
\usepackage{hyperref}

% A custom macro actually used in the body below.
\newcommand{\keyword}[1]{\textbf{#1}}
\newcommand{\Reals}{\mathbb{R}}

\title{Spectral Gaps of Random Regular Graphs}
\author{Lena Okafor \and Maren Brandt}
\date{March 2026}

\begin{document}
\maketitle

\begin{abstract}
We give an elementary proof that random $d$-regular graphs are nearly
Ramanujan with high probability, and survey the trace-method machinery
behind Friedman's theorem.
\end{abstract}

\section{Introduction}
\label{sec:intro}

Random $d$-regular graphs are \keyword{expanders} with high
probability~\cite{friedman2008}. The second eigenvalue $\lambda_2$ of the
adjacency matrix controls the mixing time of the simple random walk; see
\cite{hoory2006} for a survey. Throughout, $G = (V, E)$ denotes a graph on
$n$ vertices and $A \in \Reals^{n \times n}$ its adjacency matrix.

We make three contributions:
\begin{itemize}
  \item An \emph{elementary} second-moment argument for the spectral gap.
  \item A \texttt{numpy} implementation of the non-backtracking trace bound.
  \item Empirical evidence that the constant in \eqref{eq:gap} is tight.
\end{itemize}

\section{Main results}
\label{sec:results}

\subsection{The spectral gap}

Our main theorem refines the bound of \cite{friedman2008}.

\begin{equation}
  \label{eq:gap}
  \lambda_2(A) \le 2\sqrt{d-1} + \epsilon
\end{equation}

The proof proceeds in three steps:
\begin{enumerate}
  \item Reduce to the non-backtracking operator $B$.
  \item Bound $\operatorname{tr}(B^k)$ by counting closed walks.
  \item Apply Markov's inequality with $k \approx \log n$.
\end{enumerate}

\subsection{A two-line lemma}

For any unit vector $x \perp \mathbf{1}$ we have the variational bound
\[
  \lambda_2(A) = \max_{x \perp \mathbf{1}, \|x\| = 1} x^\top A x .
\]

The aligned chain below is standard but worth recording:
\begin{align}
  x^\top A x &= \sum_{(u,v) \in E} 2 x_u x_v \\
             &\le \sum_{(u,v) \in E} (x_u^2 + x_v^2) = d \|x\|^2 .
\end{align}

\section{Experiments}
\label{sec:experiments}

Figure~\ref{fig:gap-hist} shows the empirical distribution of
$\lambda_2 / 2\sqrt{d-1}$ over $10^4$ samples; Table~\ref{tab:constants}
lists the fitted constants. The 50\% quantile sits within 1\% of the
conjectured value.

\begin{figure}[t]
  \centering
  \includegraphics[width=0.7\linewidth]{figures/gap-histogram}
  \caption{Histogram of normalized second eigenvalues, $d = 4$.}
  \label{fig:gap-hist}
\end{figure}

\begin{table}[h]
  \centering
  \begin{tabular}{lrr}
    \toprule
    $d$ & mean & std \\
    \midrule
    3 & 0.981 & 0.012 \\
    4 & 0.987 & 0.009 \\
    \bottomrule
  \end{tabular}
  \caption{Fitted constants per degree.}
  \label{tab:constants}
\end{table}

The sampling harness is reproduced verbatim:
\begin{verbatim}
for d in 3 4 5; do
  python sample.py --degree $d --trials 10000
done
\end{verbatim}

As discussed in Section~\ref{sec:intro}, the bound \eqref{eq:gap} cannot be
improved beyond the Ramanujan threshold.

\input{appendix-proofs}

\bibliographystyle{plain}
\bibliography{refs}

\end{document}
